\documentclass[10pt,a4paper]{amsart}
\usepackage[margin=19mm]{geometry}
\usepackage{amsmath,amssymb,mathtools}
\usepackage[scr=rsfs,cal=euler]{mathalfa}
\usepackage{xfrac}
\usepackage[unicode,hidelinks,bookmarks=false]{hyperref}
\newcommand{\ie}{i.e.\ }
\newcommand{\cf}{cf.\ }
\newcommand{\resp}{respectively\ }
\newcommand{\Subsection}{\subsection}
\providecommand{\nobreakdash}{\nobreak\hbox{-}}
\setlength{\emergencystretch}{2em}
\multlinegap0pt
\usepackage{amssymb}
\usepackage{mathtools}
\newtheorem{prop}{Proposition}
\newtheorem{construction}{Construction}
\newtheorem{lem}{Lemma}
\newcommand{\s}{\sigma}
\title{On the Prym map of degree 4 cyclic covers of hyperelliptic curves}
\author{Anatoli Shatsila}
\date{Source: arXiv:2508.20838v2 (CC BY 4.0)}
\begin{document}
\maketitle

\noindent\emph{Source and licence.} On the Prym map of degree 4 cyclic covers of hyperelliptic curves. Version arXiv:2508.20838v2 (CC BY 4.0), distributed by arXiv under the Creative Commons Attribution 4.0 International licence, http://creativecommons.org/licenses/by/4.0/, as recorded in the arXiv licence metadata for this exact version and shown on the abstract page https://arxiv.org/abs/2508.20838v2. Full licence notice: \texttt{licenses/arxiv-2508.20838-CC-BY-4.0.txt}.\par\smallskip

\noindent\textit{Editorial scope.} Counted unit: the article's lem geomprym (source0.tex lines 355--368). The article's complete original proof of it is delivered with it (source0.tex lines 370--386, one \texttt{proof} environment of 17 lines). No mathematics is written, repaired or re-derived: the statement and the proof are the article's own lines, in the article's own notation, and no retained line is substituted or re-flowed.  Retained uncounted as the source-local context the proof needs: lines 237--241; lines 250--256.  One declared adaptation: exactly 1 ancillary strings inside the retained spans are omitted (image-inclusion commands and, where the source repeats a label, the repeated label definitions) and no other line differs from the source; the figure environments, with their placement, captions and labels, are kept, and the package records the omitted strings in fidelity receipts bound to the affected bodies.  No retained line contains a citation, so the wrapper carries no bibliography and no bibliography entry.  The retained lines contain the article's own diagram code, which the wrapper typesets with the standard tikz packages; no image file is delivered and the package declares no asset dependency.  Editorial formatting only, in addition: an AMS \texttt{amsart} preamble that restates the article's own declarations and macros used by the retained lines, namely the environments \texttt{prop}, \texttt{construction}, \texttt{lem} on separate counters, together with the standard packages those lines need.  The retained environments print in this document as Proposition 1, Construction 1, Lemma 1 in order of appearance, whereas the article prints the same items with the numbering of its own full document.  The complete native source of the article as distributed in the arXiv e-print for this exact version is bundled unchanged as \texttt{sources/184057/source0.tex}.
\begin{prop}
\label{bundles}
    Let $H$ be a genus $g$ hyperelliptic curve, $\{u_1, u_2, w_1, \ldots w_{2g}\}$ be the set of its Weierstrass points, $\nu = \mathcal{O}_H(u_1-u_2)$, and $k: C_{\sigma^2} \to H$ be the cover defined by $\nu$. Let $\psi: \mathbb{P}^1 \to \mathbb{P}^1$ be a double cover branched over $u_1, u_2$ and $\bar{\sigma}$ be the involution of $\mathbb{P}^1$ inducing $\psi$. With a slight abuse of notation, we denote the preimages of $w_i$ under $\psi$ by $Q_i = \{w_i, \bar\sigma(w_i)\}$. 
    Then $ \bigcup_{i=1}^{2g} Q_i$ is the set of Weierstrass points of $C_{\sigma^2}$ and, if we define $$Q := \left\{\mathcal{O}_{C_{\sigma^2}}\left(\sum_{i=1}^{2g}(-1)^iq_i\right) \ \mid \ q_i \in Q_i\right\}$$ then $$Q = \{k^*\eta \ | \ \eta^2 = \nu\}.$$
\end{prop}

\begin{construction}
\label{construction}
    Using Proposition \ref{bundles}, we can view the construction of the cover $f:C \to H$ starting from $2g+2$ Weierstrass points $W = \{u_1, u_2, w_1,\ldots, w_{2g}\}$ with two distinguished points $u_1, u_2$ and a line bundle $\xi \in Q$. We take $\nu := \mathcal{O}_H(u_1-u_2)$ and $\xi := \mathcal{O}_{C_{\sigma^2}}(\sum_{i=1}^{2g}(-1)^iq_i)$ for some $q_i \in Q_i$. For $g=2$ we have the following diagram, where ramifications are indicated above each arrow. 
    \begin{center}
        \label{diagram2}
    \end{center}
\end{construction}

\begin{lem}
\begin{enumerate}
\label{geomprym}
    \item[i)] With the notation from the diagram \ref{diagram2} we have
    $$\begin{aligned}
\ker{\mu} = \Bigg\{ 
& \left( \sum_{i \in A_+} w_i - \sum_{j \in A_-} w_j, \ \sum_{i \in A_+} \s(q_i) - \sum_{j \in A_-} \s(q_j), \ \sum_{i \in A_+} q_i - \sum_{j \in A_-} q_j \right) \\
& \in JC_{j\sigma, \sigma^2}[2] \times JC_{j\sigma^2}[2] \times JC_{j}[2] \\
& \;\Bigg|\; A_+, A_- \subset \{1,\ldots,2g\}, \ |A_+| = |A_-|, A_+ \cap A_- = \varnothing \Bigg\}.
\end{aligned}$$
    \item[ii)] Let $\pi$ be the quotient map $P \to P/K(\Xi)[2]$ and $\alpha = \pi \circ \mu$ be the composition. Then $$(P/K(\Xi)[2])^{(1,\ldots,1)} = \alpha(JC_{j}\times JC_{j\s^2})^{(1,\ldots,1)} \times (JC_{j\s, \s^2}/JC_{j\s, \s^2}[2])^{(1,\ldots,1)}.$$
\end{enumerate}
    
\end{lem}

\begin{proof}
    Recall that $h$ denotes the cover $C \to C_{\sigma^2}$. Note that $P(h) = JC_{j}^{(2,\ldots,2)} \boxplus JC_{j\sigma^2}^{(2,\ldots,2)}$ and $P(h)$ has polarization type $(2,\ldots,2)$. Therefore, it follows that $P(h) = JC_{j} \times JC_{j\sigma^2}$ with the induced polarizations and, in particular, $JC_{j} \cap JC_{j\sigma^2} = \{0\}$ in $P$.

    Since $JH$ has polarization of type $(1,4,\ldots, 4)$, it follows that $P$ has polarization of type $(1,\ldots, 1,4, \ldots, 4)$ where 4 appears $g-1$ times. Therefore, we have $\deg(\mu) = \frac{2^{g-1}\cdot 2^{g-1}\cdot 4^{g-1}}{4^{g-1}} = 4^{g-1}$. Note that $\mathbb{Z}_4^{2g-2} \simeq K(\Xi) \subset \mu(K(\mu^*\Xi))$ since $\mu$ is surjective. Since $\ker \mu \subset K(\mu^*\Xi) \simeq \mathbb{Z}_2^{2g-2} \times \mathbb{Z}_2^{2g-2} \times \mathbb{Z}_4^{2g-2}$
it follows that $\ker \mu \simeq \mathbb{Z}_2^{2g-2}$.

Denoting $\pi_j:C \to C_j, \pi_{j\sigma^2}: C \to C_{j\sigma^2}, \pi_{j\sigma, \sigma^2}: C \to C_{j\sigma, \sigma^2}$
    one readily checks that for any two disjoint subsets $A_+, A_- \in \{1,\ldots,2g\}$ of equal cardinality we have $$\pi_j^*\left(\sum_{i \in A_+} q_i - \sum_{j \in A_-} q_j\right) + \pi^*_{j\sigma^2}\left(\sum_{i \in A_+} \s(q_i) - \sum_{j \in A_-} \s(q_j)\right) + \pi_{j\sigma, \sigma^2}^*\left(\sum_{i \in A_+} w_i - \sum_{j \in A_-} w_j\right) = 0$$ so the first part of the lemma is proved. 

In order to prove the second part, we denote the corresponding polarizations on Jacobians of the decomposition of $P$ by $\Xi_{j}, \Xi_{j\sigma^2}, \Xi_{j\sigma,\sigma^2}$ and introduce the notation $$K(\Xi_{j}) = JC_{j}[2] = \langle p_1, \ldots, p_{2g-2} \rangle,$$
$$K(\Xi_{j\sigma^2}) = JC_{j\sigma^2}[2] = \langle r_1, \ldots, r_{2g-2} \rangle,$$
$$K(\Xi_{j\sigma, \sigma^2}) = JC_{j\s, \s^2}[4] = \langle v_1, \ldots, v_{2g-2} \rangle,$$
$$K(\Xi) =  \langle s_1, \ldots, s_{2g-2} \rangle.$$

Then for any $j \in \{1,\ldots, 2g-2\}$ we have $$s_j = \sum_{i=1}^{2g-2}a_i^jp_i + \sum_{i=1}^{2g-2}b_i^jr_i + \sum_{i=1}^{2g-2}c_i^jv_i$$ for some integers $a_i^j, b_i^j, c_i^j$ with $i, j \in \{1,\ldots, 2g-2\}$. In particular, $$2s_j =  2\sum_{i=1}^{2g-2}c_i^jv_i$$ Therefore, $\mathbb{Z}_2^{2g-2} \simeq K(\Xi)[2] = \langle 2s_1, \ldots, 2s_{2g-2}\rangle = \langle 2v_1, \ldots, 2v_{2g-2} \rangle = JC_{j\s, \s^2}[2]$, and $\alpha_{|JC_{j\s, \s^2}}$ is the multiplication by 2. Thus, it follows from the first part of the lemma that, up to the change of generators of $JC_{j\sigma^2}[2]$, we have  $$\ker \alpha = \langle p_1+r_1, \ldots, p_{2g-2}+r_{2g-2}, 2v_1, \ldots, 2v_{2g-2}\rangle \simeq \mathbb{Z}_2^{4g-4}.$$ This implies that $\alpha(JC_{j} \times JC_{j\s^2}) \cap \alpha(JC_{j\s,\s^2}) = \{0\}$, which finishes the proof.

\end{proof}

\end{document}
